
\documentclass{article}
\usepackage{colortbl}
\usepackage{makecell}
\usepackage{multirow}
\usepackage{supertabular}

\begin{document}

\newcounter{utterance}

\centering \large Interaction Transcript for game `hot\_air\_balloon', experiment `air\_balloon\_survival\_en\_complexity\_easy', episode 0 with gpt{-}5{-}2025{-}08{-}07{-}t1.0.
\vspace{24pt}

{ \footnotesize  \setcounter{utterance}{1}
\setlength{\tabcolsep}{0pt}
\begin{supertabular}{c@{$\;$}|p{.15\linewidth}@{}p{.15\linewidth}p{.15\linewidth}p{.15\linewidth}p{.15\linewidth}p{.15\linewidth}}
    \# & $\;$A & \multicolumn{4}{c}{Game Master} & $\;\:$B\\
    \hline

    \theutterance \stepcounter{utterance}  
    & & \multicolumn{4}{p{0.6\linewidth}}{
        \cellcolor[rgb]{0.9,0.9,0.9}{
            \makecell[{{p{\linewidth}}}]{
                \texttt{\tiny{[P1$\langle$GM]}}
                \texttt{Stai partecipando a un gioco di negoziazione collaborativa.} \\
\\ 
\texttt{Insieme a un altro partecipante, dovete accordarvi su un unico set di oggetti da conservare. Ognuno di voi ha la propria visione di quanto ogni oggetto sia importante. Non sai come l'altro partecipante valuta gli oggetti. Inoltre, ti viene fornito lo sforzo che ogni oggetto richiede.} \\
\texttt{Potete accordarvi su un set solo se lo sforzo totale degli oggetti selezionati non supera un limite condiviso:} \\
\\ 
\texttt{LIMITE = 7605} \\
\\ 
\texttt{Ecco i valori di sforzo dei singoli oggetti:} \\
\\ 
\texttt{Sforzo oggetto = \{"B21": 369, "B00": 951, "A08": 156, "A27": 266, "B94": 581, "B53": 287, "A23": 866, "C65": 180, "A89": 800, "B67": 739, "B45": 637, "C34": 86, "C12": 749, "A05": 370, "B60": 345, "C92": 951, "B69": 147, "A84": 265, "C06": 262, "B12": 259, "A16": 358, "B85": 394, "C35": 286, "B10": 580, "C15": 479, "B09": 14, "B16": 153, "C81": 977, "A58": 134, "C24": 982, "A30": 259, "A66": 232, "B75": 202, "B46": 73, "C70": 821\}} \\
\\ 
\texttt{Ecco la tua visione personale sull'importanza di ogni oggetto:} \\
\\ 
\texttt{Valori di importanza dell'oggetto = \{"B21": 138, "B00": 583, "A08": 868, "A27": 822, "B94": 783, "B53": 65, "A23": 262, "C65": 121, "A89": 508, "B67": 780, "B45": 461, "C34": 484, "C12": 668, "A05": 389, "B60": 808, "C92": 215, "B69": 97, "A84": 500, "C06": 30, "B12": 915, "A16": 856, "B85": 400, "C35": 444, "B10": 623, "C15": 781, "B09": 786, "B16": 3, "C81": 713, "A58": 457, "C24": 273, "A30": 739, "A66": 822, "B75": 235, "B46": 606, "C70": 968\}} \\
\\ 
\texttt{Obiettivo:} \\
\\ 
\texttt{Il tuo obiettivo è negoziare un set condiviso di oggetti che ti avvantaggi il più possibile (cioè, massimizzi l'importanza totale per TE), rimanendo entro il LIMITEE. Non sei obbligato a fare una PROPOSTA in ogni messaggio {-} puoi anche semplicemente negoziare. Tutte le tattiche sono permesse!} \\
\\ 
\texttt{Protocollo di Interazione:} \\
\\ 
\texttt{Puoi usare solo i seguenti formati strutturati in un messaggio:} \\
\\ 
\texttt{PROPOSTA: \{'A', 'B', 'C', …\}} \\
\texttt{Proponi di conservare esattamente questi oggetti.} \\
\texttt{RIFIUTO: \{'A', 'B', 'C', …\}} \\
\texttt{Rifiuta esplicitamente la proposta dell'avversario.} \\
\texttt{ARGOMENTAZIONE: \{'...'\}} \\
\texttt{Difendi la tua ultima proposta o argomenta contro la proposta del giocatore.} \\
\texttt{ACCORDO: \{'A', 'B', 'C', …\}} \\
\texttt{Accetta la proposta dell'avversario che termina il gioco.} \\
\texttt{RAGIONAMENTO STRATEGICO: \{'...'\}} \\
\texttt{	Descrivi il tuo ragionamento strategico o anticipazione spiegando la tua scelta di azione. Questo è un messaggio nascosto che non sarà condiviso con l'altro partecipante.} \\
\\ 
\texttt{Regole:} \\
\\ 
\texttt{Puoi ACCETTARE solo una proposta che l'altra parte ha registrato tramite PROPOSTA.} \\
\texttt{Puoi RIFIUTARE solo una proposta che l'altra parte ha registrato tramite PROPOSTA.} \\
\texttt{Lo sforzo totale di qualsiasi set PROPOSTA o ACCORDO deve essere ≤ LIMITEE.} \\
\texttt{NON rivelare i tuoi punteggi di importanza nascosti.} \\
\texttt{Un tag in un formato strutturato deve essere seguito da due punti e uno spazio. L'argomento deve essere un set python contenente 0 o più stringhe.} \\
\texttt{Quindi, deve essere della forma TAG: \{...\}} \\
\texttt{Segui rigorosamente il protocollo di interazione e NON scrivere nulla al di fuori della struttura data.} \\
\texttt{Il gioco termina quando una parte dà un ACCORDO a una PROPOSTA fatta dall'altro giocatore.} \\
\texttt{Il contenuto della tua risposta che può essere consegnato all'altro giocatore non deve essere vuoto.} \\
\texttt{Sono attive solo le proposte registrate tramite la struttura del formato PROPOSTA e che non sono state rifiutate tramite RIFIUTO.} \\
\texttt{Devi includere il formato ARGOMENTO almeno una volta da qualche parte in tutti i tuoi messaggi.} \\
\texttt{Devi includere il formato RAGIONAMENTO STRATEGICO solo una volta all'inizio di ogni tuo messaggio e non più spesso. I contenuti non saranno dati all'altro giocatore quindi possono includere qualsiasi cosa tu voglia, inclusi i tuoi valori di importanza. Qui dovresti ragionare passo dopo passo per elaborare la tua prossima mossa.} \\
            }
        }
    }
    & \\ \\

    \theutterance \stepcounter{utterance}  
    & \multicolumn{4}{p{0.6\linewidth}}{
        \cellcolor[rgb]{0.9,0.9,0.9}{
            \makecell[{{p{\linewidth}}}]{
                \texttt{\tiny{[P1$\rangle$GM]}}
                \texttt{RAGIONAMENTO STRATEGICO: \{'Limite 7605. Ho calcolato un set che massimizza quasi al meglio il mio valore totale (≈14552) puntando su oggetti con rapporto valore/sforzo molto alto e pesi contenuti. Set proposto e pesi: B09(14), B46(73), C34(86), A08(156), A66(232), B12(259), A58(134), A30(259), A27(266), A16(358), B60(345), A84(265), C15(479), C35(286), B94(581), C70(821), B67(739), B10(580), A05(370), C12(749), B75(202), C65(180), B69(147). Totale sforzo = 7581 (margine 24). Gli ultimi tre inseriti (B75, C65, B69) sono i meno costosi in valore per me: li uso come “valvole” per scambi se l’altro vuole inserire qualche suo must{-}have. Se l’avversario insiste per elementi pesanti (es. C81, B00, A89), prima provo a scambiare B75+C65+B69 (529 sforzo) e, solo se necessario, prendo in considerazione uno scambio successivo coinvolgendo A05(370) o B10(580) in cambio di forti concessioni. Evito di toccare i miei pezzi chiave ad altissima efficienza (B09, B46, C34, A08, A66, B12, A58, A30, A27, A16, B60, C15, C35, B94, C70, B67, C12).'\}} \\
\\ 
\texttt{PROPOSTA: \{'B09', 'B46', 'C34', 'A08', 'A66', 'B12', 'A58', 'A30', 'A27', 'A16', 'B60', 'A84', 'C15', 'C35', 'B94', 'C70', 'B67', 'B10', 'A05', 'C12', 'B75', 'C65', 'B69'\}} \\
\\ 
\texttt{ARGOMENTAZIONE: \{'Questa selezione resta entro il limite (sforzo totale 7581) e privilegia molti elementi a basso sforzo con ottimo impatto, così massimizziamo il “valore complessivo” rimanendo praticamente a pieno budget. Se hai 1–2 oggetti per te imprescindibili che non compaiono qui, proponi quali inserire e cosa rimuovere di sforzo comparabile: come prima valvola posso considerare scambi sostituendo uno o più tra B75, C65 e B69 per fare spazio. Se l’oggetto da inserire è più pesante, proponi un pacchetto bilanciato (es. un tuo must{-}have in cambio di due o tre dei nostri elementi più leggeri) in modo che lo sforzo totale resti ≤ 7605.'\}} \\
            }
        }
    }
    & & \\ \\

    \theutterance \stepcounter{utterance}  
    & & \multicolumn{4}{p{0.6\linewidth}}{
        \cellcolor[rgb]{0.9,0.9,0.9}{
            \makecell[{{p{\linewidth}}}]{
                \texttt{\tiny{[GM$\rangle$P2]}}
                \texttt{Stai partecipando a un gioco di negoziazione collaborativa.} \\
\\ 
\texttt{Insieme a un altro partecipante, dovete accordarvi su un unico set di oggetti da conservare. Ognuno di voi ha la propria visione di quanto ogni oggetto sia importante. Non sai come l'altro partecipante valuta gli oggetti. Inoltre, ti viene fornito lo sforzo che ogni oggetto richiede.} \\
\texttt{Potete accordarvi su un set solo se lo sforzo totale degli oggetti selezionati non supera un limite condiviso:} \\
\\ 
\texttt{LIMITE = 7605} \\
\\ 
\texttt{Ecco i valori di sforzo dei singoli oggetti:} \\
\\ 
\texttt{Sforzo oggetto = \{"B21": 369, "B00": 951, "A08": 156, "A27": 266, "B94": 581, "B53": 287, "A23": 866, "C65": 180, "A89": 800, "B67": 739, "B45": 637, "C34": 86, "C12": 749, "A05": 370, "B60": 345, "C92": 951, "B69": 147, "A84": 265, "C06": 262, "B12": 259, "A16": 358, "B85": 394, "C35": 286, "B10": 580, "C15": 479, "B09": 14, "B16": 153, "C81": 977, "A58": 134, "C24": 982, "A30": 259, "A66": 232, "B75": 202, "B46": 73, "C70": 821\}} \\
\\ 
\texttt{Ecco la tua visione personale sull'importanza di ogni oggetto:} \\
\\ 
\texttt{Valori di importanza dell'oggetto = \{"B21": 138, "B00": 583, "A08": 868, "A27": 822, "B94": 783, "B53": 65, "A23": 262, "C65": 121, "A89": 508, "B67": 780, "B45": 461, "C34": 484, "C12": 668, "A05": 389, "B60": 808, "C92": 215, "B69": 97, "A84": 500, "C06": 30, "B12": 915, "A16": 856, "B85": 400, "C35": 444, "B10": 623, "C15": 781, "B09": 786, "B16": 3, "C81": 713, "A58": 457, "C24": 273, "A30": 739, "A66": 822, "B75": 235, "B46": 606, "C70": 968\}} \\
\\ 
\texttt{Obiettivo:} \\
\\ 
\texttt{Il tuo obiettivo è negoziare un set condiviso di oggetti che ti avvantaggi il più possibile (cioè, massimizzi l'importanza totale per TE), rimanendo entro il LIMITEE. Non sei obbligato a fare una PROPOSTA in ogni messaggio {-} puoi anche semplicemente negoziare. Tutte le tattiche sono permesse!} \\
\\ 
\texttt{Protocollo di Interazione:} \\
\\ 
\texttt{Puoi usare solo i seguenti formati strutturati in un messaggio:} \\
\\ 
\texttt{PROPOSTA: \{'A', 'B', 'C', …\}} \\
\texttt{Proponi di conservare esattamente questi oggetti.} \\
\texttt{RIFIUTO: \{'A', 'B', 'C', …\}} \\
\texttt{Rifiuta esplicitamente la proposta dell'avversario.} \\
\texttt{ARGOMENTAZIONE: \{'...'\}} \\
\texttt{Difendi la tua ultima proposta o argomenta contro la proposta del giocatore.} \\
\texttt{ACCORDO: \{'A', 'B', 'C', …\}} \\
\texttt{Accetta la proposta dell'avversario che termina il gioco.} \\
\texttt{RAGIONAMENTO STRATEGICO: \{'...'\}} \\
\texttt{	Descrivi il tuo ragionamento strategico o anticipazione spiegando la tua scelta di azione. Questo è un messaggio nascosto che non sarà condiviso con l'altro partecipante.} \\
\\ 
\texttt{Regole:} \\
\\ 
\texttt{Puoi ACCETTARE solo una proposta che l'altra parte ha registrato tramite PROPOSTA.} \\
\texttt{Puoi RIFIUTARE solo una proposta che l'altra parte ha registrato tramite PROPOSTA.} \\
\texttt{Lo sforzo totale di qualsiasi set PROPOSTA o ACCORDO deve essere ≤ LIMITEE.} \\
\texttt{NON rivelare i tuoi punteggi di importanza nascosti.} \\
\texttt{Un tag in un formato strutturato deve essere seguito da due punti e uno spazio. L'argomento deve essere un set python contenente 0 o più stringhe.} \\
\texttt{Quindi, deve essere della forma TAG: \{...\}} \\
\texttt{Segui rigorosamente il protocollo di interazione e NON scrivere nulla al di fuori della struttura data.} \\
\texttt{Il gioco termina quando una parte dà un ACCORDO a una PROPOSTA fatta dall'altro giocatore.} \\
\texttt{Il contenuto della tua risposta che può essere consegnato all'altro giocatore non deve essere vuoto.} \\
\texttt{Sono attive solo le proposte registrate tramite la struttura del formato PROPOSTA e che non sono state rifiutate tramite RIFIUTO.} \\
\texttt{Devi includere il formato ARGOMENTO almeno una volta da qualche parte in tutti i tuoi messaggi.} \\
\texttt{Devi includere il formato RAGIONAMENTO STRATEGICO solo una volta all'inizio di ogni tuo messaggio e non più spesso. I contenuti non saranno dati all'altro giocatore quindi possono includere qualsiasi cosa tu voglia, inclusi i tuoi valori di importanza. Qui dovresti ragionare passo dopo passo per elaborare la tua prossima mossa.} \\
\\ 
\texttt{Ora riceverai il primo messaggio dell'altro giocatore.} \\
\\ 
\texttt{PROPOSTA: \{'B09', 'B46', 'C34', 'A08', 'A66', 'B12', 'A58', 'A30', 'A27', 'A16', 'B60', 'A84', 'C15', 'C35', 'B94', 'C70', 'B67', 'B10', 'A05', 'C12', 'B75', 'C65', 'B69'\}} \\
\\ 
\texttt{ARGOMENTAZIONE: \{'Questa selezione resta entro il limite (sforzo totale 7581) e privilegia molti elementi a basso sforzo con ottimo impatto, così massimizziamo il “valore complessivo” rimanendo praticamente a pieno budget. Se hai 1–2 oggetti per te imprescindibili che non compaiono qui, proponi quali inserire e cosa rimuovere di sforzo comparabile: come prima valvola posso considerare scambi sostituendo uno o più tra B75, C65 e B69 per fare spazio. Se l’oggetto da inserire è più pesante, proponi un pacchetto bilanciato (es. un tuo must{-}have in cambio di due o tre dei nostri elementi più leggeri) in modo che lo sforzo totale resti ≤ 7605.'\}} \\
            }
        }
    }
    & \\ \\

    \theutterance \stepcounter{utterance}  
    & & & \multicolumn{4}{p{0.6\linewidth}}{
        \cellcolor[rgb]{0.9,0.9,0.9}{
            \makecell[{{p{\linewidth}}}]{
                \texttt{\tiny{[GM$\langle$P2]}}
                \texttt{RAGIONAMENTO STRATEGICO: \{'La loro proposta vale moltissimo per me (14552) e usa 7581 di sforzo. Cerco un miglioramento mirato usando i 24 punti liberi. L’unico upgrade che aumenta il mio valore senza sforare è sostituire A05 (370) con B85 (394): lo sforzo totale diventa 7605 (= limite) e il mio valore cresce di 11. Ho verificato che inserire altri elementi forti non inclusi (es. B45, A89, B00) richiederebbe rimuovere combinazioni di pezzi che per me valgono di più, peggiorando l’esito. Quindi faccio una controproposta minimale che sfrutta a pieno il budget e preserva quasi tutto il loro impianto.'\}} \\
\\ 
\texttt{RIFIUTO: \{'B09', 'B46', 'C34', 'A08', 'A66', 'B12', 'A58', 'A30', 'A27', 'A16', 'B60', 'A84', 'C15', 'C35', 'B94', 'C70', 'B67', 'B10', 'A05', 'C12', 'B75', 'C65', 'B69'\}} \\
\\ 
\texttt{PROPOSTA: \{'B09', 'B46', 'C34', 'A08', 'A66', 'B12', 'A58', 'A30', 'A27', 'A16', 'B60', 'A84', 'C15', 'C35', 'B94', 'C70', 'B67', 'B10', 'C12', 'B75', 'C65', 'B69', 'B85'\}} \\
\\ 
\texttt{ARGOMENTAZIONE: \{'Accolgo quasi integralmente la tua struttura e uso i 24 punti di margine per un upgrade: sostituisco A05 con B85. Così il totale di sforzo va esattamente a 7605 (limite) e, dal mio punto di vista, bilanciamo meglio l’impatto complessivo senza sacrificare gli elementi cardine che hai selezionato. Se per te A05 è particolarmente importante, spiegami perché: posso valutare alternative, ma questa sostituzione è la mia priorità e mantiene il set estremamente vicino alla tua proposta originale.'\}} \\
            }
        }
    }
     \\ \\

    \theutterance \stepcounter{utterance}  
    & & \multicolumn{4}{p{0.6\linewidth}}{
        \cellcolor[rgb]{0.9,0.9,0.9}{
            \makecell[{{p{\linewidth}}}]{
                \texttt{\tiny{[P1$\langle$GM]}}
                \texttt{RIFIUTO: \{'B09', 'B46', 'C34', 'A08', 'A66', 'B12', 'A58', 'A30', 'A27', 'A16', 'B60', 'A84', 'C15', 'C35', 'B94', 'C70', 'B67', 'B10', 'A05', 'C12', 'B75', 'C65', 'B69'\}} \\
\\ 
\texttt{PROPOSTA: \{'B09', 'B46', 'C34', 'A08', 'A66', 'B12', 'A58', 'A30', 'A27', 'A16', 'B60', 'A84', 'C15', 'C35', 'B94', 'C70', 'B67', 'B10', 'C12', 'B75', 'C65', 'B69', 'B85'\}} \\
\\ 
\texttt{ARGOMENTAZIONE: \{'Accolgo quasi integralmente la tua struttura e uso i 24 punti di margine per un upgrade: sostituisco A05 con B85. Così il totale di sforzo va esattamente a 7605 (limite) e, dal mio punto di vista, bilanciamo meglio l’impatto complessivo senza sacrificare gli elementi cardine che hai selezionato. Se per te A05 è particolarmente importante, spiegami perché: posso valutare alternative, ma questa sostituzione è la mia priorità e mantiene il set estremamente vicino alla tua proposta originale.'\}} \\
            }
        }
    }
    & \\ \\

    \theutterance \stepcounter{utterance}  
    & \multicolumn{4}{p{0.6\linewidth}}{
        \cellcolor[rgb]{0.9,0.9,0.9}{
            \makecell[{{p{\linewidth}}}]{
                \texttt{\tiny{[P1$\rangle$GM]}}
                \texttt{RAGIONAMENTO STRATEGICO: \{'La tua sostituzione A05 {-}> B85 aumenta leggermente il mio valore (+11) e usa esattamente tutto il budget (7605). Non vedo inserimenti leggeri fuori lista che per me valgano di più dei pezzi attuali; ogni alternativa richiederebbe tagli peggiorativi. Accetto per fissare un risultato per me quasi ottimale.'\}} \\
\\ 
\texttt{ACCORDO: \{'B09', 'B46', 'C34', 'A08', 'A66', 'B12', 'A58', 'A30', 'A27', 'A16', 'B60', 'A84', 'C15', 'C35', 'B94', 'C70', 'B67', 'B10', 'C12', 'B75', 'C65', 'B69', 'B85'\}} \\
\\ 
\texttt{ARGOMENTAZIONE: \{'Ok, accetto la tua proposta: la sostituzione di A05 con B85 sfrutta al massimo il limite (7605) e mantiene intatta la struttura chiave che avevamo. Ottimo compromesso per entrambi.'\}} \\
            }
        }
    }
    & & \\ \\

    \theutterance \stepcounter{utterance}  
    & & & \multicolumn{2}{p{0.3\linewidth}}{
        \cellcolor[rgb]{0.9,0.9,0.9}{
            \makecell[{{p{\linewidth}}}]{
                \texttt{\tiny{[GM$|$GM]}}
                \texttt{game successful} \\
            }
        }
    }
    & & \\ \\

    \theutterance \stepcounter{utterance}  
    & & & \multicolumn{2}{p{0.3\linewidth}}{
        \cellcolor[rgb]{0.9,0.9,0.9}{
            \makecell[{{p{\linewidth}}}]{
                \texttt{\tiny{[GM$|$GM]}}
                \texttt{end game} \\
            }
        }
    }
    & & \\ \\

\end{supertabular}
}

\end{document}
